% ----------------------------------------------------------
% Introdução (exemplo de capítulo sem numeração, mas presente no Sumário)
% ----------------------------------------------------------
\chapter{Introdução}
% ----------------------------------------------------------

O capítulo de introdução é responsável por contextualizar o assunto a ser abordado e convencer o leitor da importância do trabalho.
Ele é composto por um texto inicial e por seções que tratam o problema, sua relevância e justificativa, os objetivos do trabalho e, normalmente, a organização da monografia.
Esse capítulo, na sua totalidade, tem função de vender o seu trabalho para o leitor e convencê-lo de sua importância.

O texto inicial do capítulo fica a critério do autor, devendo fornecer uma visão geral do trabalho.
Porém, as boas práticas acadêmicas sugerem o uso de uma estrutura semelhante a uma pirâmide invertida.
Assim, inicia-se com um assunto amplo, mas relacionado ao trabalho.
Em seguida, ocorre um afunilamento em direção ao tema principal até chegar ao problema que será abordado.
Neste ponto é feita uma pequena explicação do problema, destacando sua importância e quais as soluções existentes até o momento (se existirem), fazendo citações de trabalhos que tratam de assuntos similares.
O próximo passo é apresentar o seu trabalho com uma explicação breve do que será discutido e mostrar o que tem de diferente em relação aos demais trabalhos que tratam de assuntos similares.

Não existe um limite mínimo nem máximo para esta parte do capítulo de introdução, porém, um tamanho recomendado para TCC é de aproximadamente duas páginas.
Em geral, esse tamanho é o suficiente para contextualizar e mostrar a importância do trabalho. 

\section{Apresentação do problema}

Essa seção apresenta o problema de forma mais detalhada, lembrando que o leitor já teve um primeiro contato no texto inicial do capítulo.
O foco dessa seção é o problema e não sua importância e nem sua solução.
Esse texto é importante para que o leitor entenda o tamanho do problema.
No caso do membro da banca que irá avaliar o trabalho, essa seção é fundamental para ele entender o tamanho que o problema abordado pelo trabalho tem.
Normalmente, de dois a três parágrafos é suficiente para esta seção, porém, não existe limitação técnica e o autor tem total liberdade.

\section{Justificativa e relevância}

Nessa seção o foco muda, sai do problema e passa para o trabalho.
Aqui é discutido o motivo de estudar o tema apresentado e sua relevância.
O leitor deve finalizar essa seção convencido de que o trabalho posui contribuições e de qual é o nível dessa contribuição.
Do ponto de vista do membro da banca que irá avaliar o trabalho, essa seção é fundamental para formar o conceito nele do quanto o seu trabalho contribui.
Similar à apresentação do problema, normalmente, dois a três parágrafos são o suficiente.
Porém, o autor tem liberdade de escrever o quanto achar necessário para justificar seu trabalho.

\section{Objetivos}

Nessa seção são apresentados os objetivos do trabalho, sendo divididos em objetivos geral e específicos.
A definição dos objetivos do estudo deve ser clara e concisa.
Um objetivo modelo é aquele que informa em poucas frases precisamente aquilo que se pretende estudar.
Os enunciados devem ser apresentados com verbos no infinitivo indicando a ação a ser efetuada, exemplo: verificar, elaborar, definir, relacionar, demonstrar, identificar, desenvolver, formular, modelar, diagnosticar, determinar, definir, descrever, avaliar.

\subsection{Objetivo Geral}

Esta subseção descreve de forma global o objetivo a ser alcançado.
Ele deve iniciar com o verbo no infinitivo e deve estar diretamente relacionado ao tema do trabalho.

\subsection{Objetivos específicos}

Esta subseção lista os objetivos específicos do trabalho.
Pode ser em formato de lista ou em texto corrido, porém, cada objetivo deve iniciar com o verbo no infinitivo.
As boas práticas acadêmicas sugerem, para um TCC, o uso de dois, no máximo, três objetivos específicos.
Aqui deve-se tomar cuidado para não confundir o que necessita ser feito para atingir um objetivo com um objetivo específico.
Por exemplo, fazer uma revisão da literatura sobre um assunto específico pode não ser um objetivo do seu trabalho.
Pode ser apenas um meio de preparar o leitor para o real objetivo.
Normalmente, os objetivos específicos são tratados no Capítulo que detalha os resultados.

Se for usar lista para os objetivos específicos, a lista pode ser do tipo:
\begin{itemize}
    \item Primeiro objetivo específico
    \item Segundo objetivo específico
\end{itemize}

\section{Etapas do trabalho (ou organização da Monografia) (opcional)}

Esta seção é opcional e é utilizada para indicar ao leitor como o trabalho está organizado.
Normalmente, é realizada uma pequena descrição de cada capítulo do trabalho.
Cuidado para a descrição de cada capítulo não ficar muito pequena nem muito grande.